% Página 52 do PDF de Gabriel
\begin{frame}[plain,noframenumbering,label=original-52]
  \begin{tikzpicture}[remember picture,overlay,x=1bp,y=-1bp]
  \begin{scope}[shift={(current page.north west)}]
  \Fundo{52}
  \Texto{44.456}{41.894}{123.698}{8.518}{b}{34E0C8}{PROBABILIDADE E MODELOS}
  \Texto{43.039}{116.125}{63.147}{56.787}{b}{FFB014}{07}
  \Texto{43.039}{154.533}{162.238}{24.608}{b}{FFFFFF}{Probabilidade}
  \Texto{43.039}{234.662}{46.285}{7.572}{b}{34E0C8}{EMASC 2026}
  \Texto{91.429}{234.662}{2.65}{7.572}{b}{FFB014}{\textbullet{}}
  \Texto{96.176}{234.662}{21.45}{7.572}{b}{34E0C8}{LAME}
  \Link{42.042}{135.586}{164.229}{20.251}{original-52}
  \end{scope}
  \end{tikzpicture}
\end{frame}


% Página 53 do PDF de Gabriel
\begin{frame}[plain,noframenumbering,label=original-53]
  \begin{tikzpicture}[remember picture,overlay,x=1bp,y=-1bp]
  \begin{scope}[shift={(current page.north west)}]
  \Fundo{53}
  \Texto{40.175}{27.942}{178.584}{17.983}{b}{2C2084}{Conjuntos e eventos}
  \Texto{35.433}{63.506}{256.256}{8.518}{i}{737696}{Probabilidade fala de eventos; eventos são conjuntos de resultados.}
  \Texto{73.648}{116.403}{5.682}{8.518}{i}{262848}{A}
  \Texto{138.655}{116.403}{5.682}{8.518}{i}{262848}{B}
  \Texto{37.21}{247.463}{40.499}{6.625}{b}{737696}{EMASC 2026}
  \Texto{80.653}{247.463}{2.319}{6.625}{r}{FFB014}{\textbullet{}}
  \Texto{85.913}{247.463}{117.821}{6.625}{r}{737696}{Liga Acadêmica Matemática \& Empresa}
  \Texto{206.682}{247.463}{2.319}{6.625}{r}{FFB014}{\textbullet{}}
  \Texto{211.943}{247.463}{13.25}{6.625}{r}{737696}{UFS}
  \Texto{397.549}{248.132}{7.367}{6.625}{b}{2C2084}{30}
  \Texto{406.02}{248.132}{10.315}{6.625}{r}{737696}{/ 40}
  \end{scope}
  \end{tikzpicture}
\end{frame}


% Página 54 do PDF de Gabriel
\begin{frame}[plain,noframenumbering,label=original-54]
  \begin{tikzpicture}[remember picture,overlay,x=1bp,y=-1bp]
  \begin{scope}[shift={(current page.north west)}]
  \Fundo{54}
  \Texto{40.175}{27.803}{87.916}{17.983}{b}{2C2084}{Contagem}
  \Texto{35.433}{63.506}{310.124}{8.518}{i}{737696}{Quando todos os resultados são igualmente prováveis, probabilidade é contagem.}
  \Texto{210.648}{100.978}{66.564}{9.465}{r}{262848}{casos possíveis}
  \Texto{193.725}{163.384}{4.732}{9.465}{i}{262848}{k}
  \Texto{35.433}{186.115}{65.21}{9.465}{b}{262848}{Combinações:}
  \Texto{103.898}{186.115}{164.37}{9.465}{r}{262848}{escolher 2 alunos de uma turma de 20:}
  \Texto{277.929}{189.55}{3.684}{6.625}{r}{262848}{2}
  \Texto{49.227}{209.024}{49.711}{8.518}{r}{000000}{O coeficiente}
  \Texto{105.529}{212.067}{2.839}{5.679}{i}{000000}{k}
  \Texto{115.53}{209.024}{49.285}{8.518}{r}{000000}{conta grupos}
  \Texto{167.183}{209.024}{45.273}{8.518}{b}{000000}{sem ordem}
  \Texto{212.456}{209.024}{150.505}{8.518}{r}{000000}{; ele reaparece na distribuição binomial.}
  \Texto{37.21}{247.463}{40.499}{6.625}{b}{737696}{EMASC 2026}
  \Texto{80.653}{247.463}{2.319}{6.625}{r}{FFB014}{\textbullet{}}
  \Texto{85.913}{247.463}{117.821}{6.625}{r}{737696}{Liga Acadêmica Matemática \& Empresa}
  \Texto{206.682}{247.463}{2.319}{6.625}{r}{FFB014}{\textbullet{}}
  \Texto{211.943}{247.463}{13.25}{6.625}{r}{737696}{UFS}
  \Texto{397.549}{248.132}{7.367}{6.625}{b}{2C2084}{31}
  \Texto{406.02}{248.132}{10.315}{6.625}{r}{737696}{/ 40}
  \end{scope}
  \end{tikzpicture}
\end{frame}


% Página 55 do PDF de Gabriel
\begin{frame}[plain,noframenumbering,label=original-55]
  \begin{tikzpicture}[remember picture,overlay,x=1bp,y=-1bp]
  \begin{scope}[shift={(current page.north west)}]
  \Fundo{55}
  \Texto{40.175}{27.803}{222.48}{17.983}{b}{2C2084}{Probabilidade: definições}
  \Texto{241.22}{105.461}{5.262}{9.465}{i}{262848}{n}
  \Texto{49.227}{156.029}{357.083}{8.518}{r}{000000}{A frequência relativa estima a probabilidade de um evento. A probabilidade condicional res-}
  \Texto{49.227}{166.988}{133.792}{8.518}{r}{000000}{tringe o espaço amostral ao evento}
  \Texto{185.623}{166.988}{5.682}{8.518}{i}{000000}{B}
  \Texto{191.683}{166.988}{214.628}{8.518}{r}{000000}{; independência significa que essa restrição não altera a}
  \Texto{49.227}{177.947}{39.302}{8.518}{r}{000000}{chance de}
  \Texto{90.896}{177.947}{5.682}{8.518}{i}{000000}{A}
  \Texto{96.578}{177.947}{2.368}{8.518}{r}{000000}{.}
  \Texto{37.21}{247.463}{40.499}{6.625}{b}{737696}{EMASC 2026}
  \Texto{80.653}{247.463}{2.319}{6.625}{r}{FFB014}{\textbullet{}}
  \Texto{85.913}{247.463}{117.821}{6.625}{r}{737696}{Liga Acadêmica Matemática \& Empresa}
  \Texto{206.682}{247.463}{2.319}{6.625}{r}{FFB014}{\textbullet{}}
  \Texto{211.943}{247.463}{13.25}{6.625}{r}{737696}{UFS}
  \Texto{397.549}{248.132}{7.367}{6.625}{b}{2C2084}{32}
  \Texto{406.02}{248.132}{10.315}{6.625}{r}{737696}{/ 40}
  \end{scope}
  \end{tikzpicture}
\end{frame}


% Página 56 do PDF de Gabriel
\begin{frame}[plain,noframenumbering,label=original-56]
  \begin{tikzpicture}[remember picture,overlay,x=1bp,y=-1bp]
  \begin{scope}[shift={(current page.north west)}]
  \Fundo{56}
  \Texto{40.175}{29.691}{224.476}{17.983}{b}{2C2084}{Probabilidade condicional}
  \Texto{35.433}{63.506}{232.533}{8.518}{i}{737696}{Saber que um evento ocorreu pode alterar a chance de outro.}
  \Texto{49.227}{203.725}{357.083}{8.518}{r}{000000}{Entre os estudantes que usam mais de 6 h de redes, a proporção com pontuação de saúde}
  \Texto{49.227}{225.643}{60.112}{8.518}{b}{000000}{independentes}
  \Texto{109.338}{225.643}{2.368}{8.518}{r}{000000}{.}
  \Texto{37.21}{247.463}{40.499}{6.625}{b}{737696}{EMASC 2026}
  \Texto{80.653}{247.463}{2.319}{6.625}{r}{FFB014}{\textbullet{}}
  \Texto{85.913}{247.463}{117.821}{6.625}{r}{737696}{Liga Acadêmica Matemática \& Empresa}
  \Texto{206.682}{247.463}{2.319}{6.625}{r}{FFB014}{\textbullet{}}
  \Texto{211.943}{247.463}{13.25}{6.625}{r}{737696}{UFS}
  \Texto{397.549}{248.132}{7.367}{6.625}{b}{2C2084}{33}
  \Texto{406.02}{248.132}{10.315}{6.625}{r}{737696}{/ 40}
  \end{scope}
  \end{tikzpicture}
\end{frame}


% Páginas 57 a 58 do PDF de Gabriel
\begin{frame}[plain,noframenumbering,label=original-57]
  \begin{tikzpicture}[remember picture,overlay,x=1bp,y=-1bp]
  \begin{scope}[shift={(current page.north west)}]
  \only<1>{\Fundo{57}}
  \only<2>{\Fundo{58}}
  \Texto{40.175}{27.803}{178.512}{17.983}{b}{2C2084}{Binomial: simulação}
  \Texto{35.433}{63.506}{252.337}{8.518}{i}{737696}{Em uma turma de 20, quantos alunos esperamos com sono curto?}
  \only<1>{\Texto{49.227}{209.198}{357.083}{8.518}{r}{000000}{Turmas de 20 estudantes sorteadas dos dados; conta-se quantos dormem menos de 6 h,}}
  \only<1>{\Texto{49.227}{220.157}{190.198}{8.518}{r}{000000}{evento que ocorre em cerca de 33\% dos registros.}}
  \Texto{37.21}{247.463}{40.499}{6.625}{b}{737696}{EMASC 2026}
  \Texto{80.653}{247.463}{2.319}{6.625}{r}{FFB014}{\textbullet{}}
  \Texto{85.913}{247.463}{117.821}{6.625}{r}{737696}{Liga Acadêmica Matemática \& Empresa}
  \Texto{206.682}{247.463}{2.319}{6.625}{r}{FFB014}{\textbullet{}}
  \Texto{211.943}{247.463}{13.25}{6.625}{r}{737696}{UFS}
  \Texto{397.549}{248.132}{7.367}{6.625}{b}{2C2084}{34}
  \Texto{406.02}{248.132}{10.315}{6.625}{r}{737696}{/ 40}
  \only<2>{\Texto{49.227}{209.394}{357.083}{8.518}{r}{000000}{Repetindo o sorteio 10.000 vezes, as frequências observadas aproximam-se das probabilida-}}
  \only<2>{\Texto{49.227}{220.353}{25.571}{8.518}{r}{000000}{des da}}
  \only<2>{\Texto{77.165}{220.353}{85.504}{8.518}{b}{000000}{distribuição binomial}}
  \only<2>{\Texto{162.668}{220.353}{2.368}{8.518}{r}{000000}{.}}
  \end{scope}
  \end{tikzpicture}
\end{frame}


% Página 59 do PDF de Gabriel
\begin{frame}[plain,noframenumbering,label=original-59]
  \begin{tikzpicture}[remember picture,overlay,x=1bp,y=-1bp]
  \begin{scope}[shift={(current page.north west)}]
  \Fundo{59}
  \Texto{40.175}{29.916}{156.502}{17.983}{b}{2C2084}{Binomial: fórmula}
  \Texto{109.687}{80.832}{4.732}{9.465}{i}{262848}{k}
  \Texto{236.631}{87.561}{55.553}{6.974}{t}{000000}{n,\ p\ =\ 20,\ 0.33}
  \Texto{236.631}{105.494}{88.881}{6.974}{t}{000000}{stats.binom.pmf(0,\ n,\ p)}
  \Texto{336.637}{105.494}{29.625}{6.974}{c}{737696}{\#\ P(X=0)}
  \Texto{236.631}{114.46}{85.178}{6.974}{t}{000000}{stats.binom.sf(9,\ n,\ p)}
  \Texto{336.637}{114.46}{37.031}{6.974}{c}{737696}{\#\ P(X>=10)}
  \Texto{236.631}{123.427}{51.843}{6.974}{t}{000000}{n*p,\ n*p*(1-p)}
  \Texto{336.637}{123.427}{66.663}{6.974}{c}{737696}{\#\ media,\ variancia}
  \Texto{49.227}{169.395}{32.952}{8.966}{t}{2C2084}{sf(k-1)}
  \Texto{82.17}{169.395}{103.306}{8.518}{r}{000000}{. O modelo exige tentativas}
  \Texto{187.843}{169.395}{60.112}{8.518}{b}{000000}{independentes}
  \Texto{250.322}{169.395}{58.238}{8.518}{r}{000000}{e probabilidade}
  \Texto{310.927}{169.395}{40.231}{8.518}{b}{000000}{constante}
  \Texto{351.157}{169.395}{2.368}{8.518}{r}{000000}{.}
  \Texto{37.21}{247.463}{40.499}{6.625}{b}{737696}{EMASC 2026}
  \Texto{80.653}{247.463}{2.319}{6.625}{r}{FFB014}{\textbullet{}}
  \Texto{85.913}{247.463}{117.821}{6.625}{r}{737696}{Liga Acadêmica Matemática \& Empresa}
  \Texto{206.682}{247.463}{2.319}{6.625}{r}{FFB014}{\textbullet{}}
  \Texto{211.943}{247.463}{13.25}{6.625}{r}{737696}{UFS}
  \Texto{397.549}{248.132}{7.367}{6.625}{b}{2C2084}{35}
  \Texto{406.02}{248.132}{10.315}{6.625}{r}{737696}{/ 40}
  \end{scope}
  \end{tikzpicture}
\end{frame}


% Página 60 do PDF de Gabriel
\begin{frame}[plain,noframenumbering,label=original-60]
  \begin{tikzpicture}[remember picture,overlay,x=1bp,y=-1bp]
  \begin{scope}[shift={(current page.north west)}]
  \Fundo{60}
  \Texto{40.175}{27.803}{171.805}{17.983}{b}{2C2084}{Poisson: simulação}
  \Texto{35.433}{63.506}{275.421}{8.518}{i}{737696}{Para eventos raros em grupos grandes, contamos ocorrências por grupo.}
  \Texto{49.227}{203.827}{298.454}{8.518}{r}{000000}{Grupos de 100 estudantes; conta-se quantos dormem menos de 4 h, em média}
  \Texto{349.87}{203.827}{54.072}{8.518}{b}{000000}{2,3 por grupo}
  \Texto{403.941}{203.827}{2.368}{8.518}{r}{000000}{.}
  \Texto{49.227}{214.786}{260.429}{8.518}{r}{000000}{A curva corresponde à distribuição de Poisson com essa taxa média.}
  \Texto{37.21}{247.463}{40.499}{6.625}{b}{737696}{EMASC 2026}
  \Texto{80.653}{247.463}{2.319}{6.625}{r}{FFB014}{\textbullet{}}
  \Texto{85.913}{247.463}{117.821}{6.625}{r}{737696}{Liga Acadêmica Matemática \& Empresa}
  \Texto{206.682}{247.463}{2.319}{6.625}{r}{FFB014}{\textbullet{}}
  \Texto{211.943}{247.463}{13.25}{6.625}{r}{737696}{UFS}
  \Texto{397.549}{248.132}{7.367}{6.625}{b}{2C2084}{36}
  \Texto{406.02}{248.132}{10.315}{6.625}{r}{737696}{/ 40}
  \end{scope}
  \end{tikzpicture}
\end{frame}


% Página 61 do PDF de Gabriel
\begin{frame}[plain,noframenumbering,label=original-61]
  \begin{tikzpicture}[remember picture,overlay,x=1bp,y=-1bp]
  \begin{scope}[shift={(current page.north west)}]
  \Fundo{61}
  \Texto{40.175}{29.916}{149.794}{17.983}{b}{2C2084}{Poisson: fórmula}
  \Texto{236.631}{87.561}{33.328}{6.974}{t}{000000}{lam\ =\ 2.3}
  \Texto{236.631}{105.494}{92.584}{6.974}{t}{000000}{stats.poisson.pmf(0,\ lam)}
  \Texto{340.341}{105.494}{59.256}{6.974}{c}{737696}{\#\ P(X=0)\ \textasciitilde{}\ 0.100}
  \Texto{236.631}{114.46}{88.881}{6.974}{t}{000000}{stats.poisson.sf(4,\ lam)}
  \Texto{340.341}{114.46}{62.96}{6.974}{c}{737696}{\#\ P(X>=5)\ \textasciitilde{}\ 0.084}
  \Texto{49.227}{149.01}{217.541}{8.518}{r}{000000}{A distribuição de Poisson é o limite da binomial quando}
  \Texto{270.164}{149.01}{4.736}{8.518}{i}{000000}{n}
  \Texto{278.434}{149.01}{27.326}{8.518}{r}{000000}{cresce,}
  \Texto{309.409}{149.01}{4.736}{8.518}{i}{000000}{p}
  \Texto{317.777}{149.01}{75.427}{8.518}{r}{000000}{diminui e o produto}
  \Texto{396.599}{149.01}{9.472}{8.518}{i}{000000}{np}
  \Texto{49.227}{159.969}{309.289}{8.518}{r}{000000}{permanece constante. Nas contagens por grupo, média e variância são próximas.}
  \Texto{37.21}{247.463}{40.499}{6.625}{b}{737696}{EMASC 2026}
  \Texto{80.653}{247.463}{2.319}{6.625}{r}{FFB014}{\textbullet{}}
  \Texto{85.913}{247.463}{117.821}{6.625}{r}{737696}{Liga Acadêmica Matemática \& Empresa}
  \Texto{206.682}{247.463}{2.319}{6.625}{r}{FFB014}{\textbullet{}}
  \Texto{211.943}{247.463}{13.25}{6.625}{r}{737696}{UFS}
  \Texto{397.549}{248.132}{7.367}{6.625}{b}{2C2084}{37}
  \Texto{406.02}{248.132}{10.315}{6.625}{r}{737696}{/ 40}
  \end{scope}
  \end{tikzpicture}
\end{frame}


% Página 62 do PDF de Gabriel
\begin{frame}[plain,noframenumbering,label=original-62]
  \begin{tikzpicture}[remember picture,overlay,x=1bp,y=-1bp]
  \begin{scope}[shift={(current page.north west)}]
  \Fundo{62}
  \Texto{40.175}{29.691}{179.142}{17.983}{b}{2C2084}{Binomial ou Poisson}
  \Texto{35.433}{63.506}{245.83}{8.518}{i}{737696}{Dois modelos de contagem: a escolha depende do que se conta.}
  \Texto{38.534}{100.643}{39.969}{9.465}{b}{027499}{Binomial}
  \Texto{38.534}{115.586}{25.251}{9.465}{r}{000000}{Conta}
  \Texto{66.416}{115.586}{59.447}{9.465}{b}{000000}{sucessos em}
  \Texto{128.493}{115.586}{5.262}{9.465}{i}{000000}{n}
  \Texto{136.542}{115.586}{43.991}{9.465}{b}{000000}{tentativas}
  \Texto{180.532}{115.586}{2.631}{9.465}{r}{000000}{.}
  \Texto{38.534}{127.542}{65.116}{9.465}{r}{000000}{Limite superior:}
  \Texto{106.904}{127.542}{5.262}{9.465}{i}{000000}{n}
  \Texto{112.322}{127.542}{2.631}{9.465}{r}{000000}{.}
  \Texto{38.534}{171.377}{158.767}{8.518}{r}{737696}{Ex.: estudantes com sono curto em turma}
  \Texto{38.534}{183.332}{23.68}{8.518}{r}{737696}{de 20.}
  \Texto{250.943}{100.642}{36.438}{9.465}{b}{CC8D10}{Poisson}
  \Texto{250.943}{115.586}{25.251}{9.465}{r}{000000}{Conta}
  \Texto{278.825}{115.586}{128.774}{9.465}{b}{000000}{ocorrências em um intervalo}
  \Texto{407.597}{115.586}{2.631}{9.465}{r}{000000}{.}
  \Texto{250.943}{127.542}{83.288}{9.465}{r}{000000}{Sem limite superior.}
  \Texto{250.943}{159.422}{154.253}{8.518}{r}{737696}{Ex.: chamadas por hora em uma central.}
  \Texto{49.227}{213.169}{90.904}{8.518}{r}{000000}{Ambas supõem eventos}
  \Texto{142.498}{213.169}{60.112}{8.518}{b}{000000}{independentes}
  \Texto{204.977}{213.169}{64.745}{8.518}{r}{000000}{e taxa constante.}
  \Texto{37.21}{247.463}{40.499}{6.625}{b}{737696}{EMASC 2026}
  \Texto{80.653}{247.463}{2.319}{6.625}{r}{FFB014}{\textbullet{}}
  \Texto{85.913}{247.463}{117.821}{6.625}{r}{737696}{Liga Acadêmica Matemática \& Empresa}
  \Texto{206.682}{247.463}{2.319}{6.625}{r}{FFB014}{\textbullet{}}
  \Texto{211.943}{247.463}{13.25}{6.625}{r}{737696}{UFS}
  \Texto{397.549}{248.132}{7.367}{6.625}{b}{2C2084}{38}
  \Texto{406.02}{248.132}{10.315}{6.625}{r}{737696}{/ 40}
  \end{scope}
  \end{tikzpicture}
\end{frame}
